%=========== ΛΥΣΗ - 25ο ΘΕΜΑ Δ ===========
\begin{thema}{Δ}
\begin{erwthma}
\item Για $x=0$ η δοσμένη διπλή ανισότητα δίνει
\[
1\leq f(0)\leq1,
\]
οπότε
\[
{f(0)=1}.
\]
Επίσης, από τη σχέση
\[
f'(x)-f(x)=2xe^x
\]
για $x=0$ παίρνουμε
\[
f'(0)-f(0)=0.
\]
Άρα
\[
{f'(0)=1}.
\]

Για κάθε $x\in\mathbb{R}$ έχουμε
\begin{align*}
\left(f(x)e^{-x}\right)'
&=f'(x)e^{-x}-f(x)e^{-x}\\
&=e^{-x}\left(f'(x)-f(x)\right)\\
&=2x.
\end{align*}
Επειδή
\[
2x=(x^2)',
\]
σύμφωνα με τις συνέπειες του Θεωρήματος Μέσης Τιμής, υπάρχει σταθερά $c\in\mathbb{R}$ τέτοια ώστε
\[
f(x)e^{-x}=x^2+c.
\]
Για $x=0$ και $f(0)=1$ έχουμε $c=1$. Επομένως
\[
{f(x)=(x^2+1)e^x},
\quad x\in\mathbb{R}.
\]

\item Για κάθε $x\in\mathbb{R}$ είναι
\begin{align*}
f'(x)
&=\left[(x^2+1)e^x\right]'\\
&=2xe^x+(x^2+1)e^x\\
&=e^x(x+1)^2.
\end{align*}
Άρα
\[
f'(x)\geq0
\]
για κάθε $x\in\mathbb{R}$, με την ισότητα να ισχύει μόνο για $x=-1$. Επομένως η $f$ είναι \textbf{γνησίως αύξουσα} στο $\mathbb{R}$.

Επειδή $f(0)=1$, από τη γνήσια μονοτονία της $f$ έχουμε
\[
f(u)>1\Leftrightarrow u>0.
\]
Άρα
\begin{align*}
f\left(f(x^2-2)-1\right)>1
&\Leftrightarrow f(x^2-2)-1>0\\
&\Leftrightarrow f(x^2-2)>1\\
&\Leftrightarrow x^2-2>0\\
&\Leftrightarrow x<-\sqrt{2}\ \text{ή}\ x>\sqrt{2}.
\end{align*}
Επομένως το σύνολο λύσεων είναι
\[
{x\in(-\infty,-\sqrt{2})\cup(\sqrt{2},+\infty)}.
\]

\item Για κάθε $x\in\mathbb{R}$ είναι
\begin{align*}
f''(x)
&=\left[e^x(x+1)^2\right]'\\
&=e^x(x+1)^2+2e^x(x+1)\\
&=e^x(x+1)(x+3).
\end{align*}
Για $x>0$ έχουμε
\[
f''(x)>0.
\]
Άρα η $f$ είναι αυστηρά κυρτή στο $(0,+\infty)$.

Επειδή το $x+1$ είναι το μέσο των $x$ και $x+2$, από την αυστηρή κυρτότητα της $f$ έχουμε
\[
f(x+1)<\frac{f(x)+f(x+2)}{2}.
\]
Επομένως, για κάθε $x>0$,
\[
{f(x+2)+f(x)>2f(x+1)}.
\]

\item Από τον τύπο της $f$ έχουμε
\begin{align*}
I
&=\int_0^1(x^2+1)e^x\d x\\
&=\left[(x^2-2x+3)e^x\right]_0^1\\
&=2e-3.
\end{align*}
Άρα
\[
{I=2e-3}.
\]
\end{erwthma}
\end{thema}
%=========== ΤΕΛΟΣ ΛΥΣΗΣ - 25ο ΘΕΜΑ Δ ===========
